\section{Proofs for Sections~\ref{sec:model}--\ref{sec:identification}}\label{app:proofs-measurement}

Throughout, $\operatorname{offdiag}(\bm M)$ denotes the off-diagonal entries of a square matrix $\bm M$, and two matrices are \emph{off-diagonally equal} if their off-diagonal entries agree. For a domain $\mathcal D\in\{\mathcal S,\mathcal I\}$ we write $\bm w=\bm\lambda_{\mathcal DG}$ and $\bm s=\bm\lambda_{\mathcal Ds}$ for the standardised general and specific loadings of its items, and we drop the superscript $\mathcal D$ when the domain is clear. Response patterns take finitely many values, so every expectation below is a finite sum. All models considered are $C^\infty$ in their parameters, because $F$ is smooth and the latent distributions are Gaussian.

\subsection{Algebraic lemmas}\label{app:pm-lemmas}

\begin{lemma}[Rank-one factorisations]\label{lem:app-rank1}
If $\bm u\bm v^\top=\tilde{\bm u}\tilde{\bm v}^\top\neq\bm0$, then $\tilde{\bm u}=t\bm u$ and $\tilde{\bm v}=\bm v/t$ for a unique $t\neq0$.
\end{lemma}

\begin{proof}
Choose $k,l$ with $u_kv_l\neq0$; then $\tilde u_k\tilde v_l=u_kv_l\neq0$. For every $k'$, $\tilde u_{k'}=\tilde u_{k'}\tilde v_l/\tilde v_l=u_{k'}v_l/\tilde v_l$, so $\tilde{\bm u}=t\bm u$ with $t=v_l/\tilde v_l$. Likewise $\tilde{\bm v}=(u_k/\tilde u_k)\bm v=\bm v/t$. Uniqueness is clear because $\bm u\neq\bm0$.
\end{proof}

\begin{lemma}[One-factor off-diagonal structure]\label{lem:app-onefactor}
Let $\bm u\in\mathbb R^p$ have at least three non-zero entries. If $\tilde{\bm u}\tilde{\bm u}^\top$ and $\bm u\bm u^\top$ are off-diagonally equal, then $\tilde{\bm u}=\pm\bm u$.
\end{lemma}

\begin{proof}
Let $N=\{k:u_k\neq0\}$, so $|N|\ge3$. For $k\in N$ pick distinct $l,m\in N\setminus\{k\}$. Then $\tilde u_l\tilde u_m=u_lu_m\neq0$ and
$\tilde u_k^2=(\tilde u_k\tilde u_l)(\tilde u_k\tilde u_m)/(\tilde u_l\tilde u_m)=u_k^2$. For $k\notin N$ and any $l\in N$, $\tilde u_k\tilde u_l=0$ with $\tilde u_l\neq0$, so $\tilde u_k=0$. Finally, $\tilde u_k\tilde u_l=u_ku_l$ for $k,l\in N$ gives $\operatorname{sgn}(\tilde u_k/u_k)=\operatorname{sgn}(\tilde u_l/u_l)$, so the sign change is common to all entries.
\end{proof}

\begin{lemma}[Tetrads]\label{lem:app-tetrad}
For $\bm w,\bm s\in\mathbb R^p$, $\delta\in\mathbb R$ and $\bm M=\bm s\bm s^\top+\delta\bm w\bm w^\top$, define $D_{kl}=w_ks_l-s_kw_l$. For distinct $k_1,\dots,k_4$,
\[
  M_{k_1k_2}M_{k_3k_4}-M_{k_1k_3}M_{k_2k_4}=\delta\,D_{k_1k_4}D_{k_2k_3},
\]
whereas the same tetrad of any rank-one matrix $\bm u\bm u^\top$ vanishes.
\end{lemma}

\begin{proof}
Write $1,\dots,4$ for $k_1,\dots,k_4$. Expanding, the terms in $\delta^0$ ($s_1s_2s_3s_4$) and in $\delta^2$ ($w_1w_2w_3w_4$) cancel, leaving
$\delta\,(s_1s_2w_3w_4+w_1w_2s_3s_4-s_1s_3w_2w_4-w_1w_3s_2s_4)$. On the other hand,
$D_{14}D_{23}=(w_1s_4-s_1w_4)(w_2s_3-s_2w_3)=w_1w_2s_3s_4-w_1w_3s_2s_4-s_1s_3w_2w_4+s_1s_2w_3w_4$, which is the same bracket. For $\bm u\bm u^\top$ the tetrad is $u_1u_2u_3u_4-u_1u_3u_2u_4=0$.
\end{proof}

\begin{lemma}[When heterogeneity fails]\label{lem:app-classes}
Let a domain contain $p\ge4$ items whose loading vectors $(w_k,s_k)$ are all non-zero. There exist distinct $k_1,\dots,k_4$ with $D_{k_1k_4}D_{k_2k_3}\neq0$ if and only if it is not the case that all items, with at most one exception, lie on one line through the origin.
\end{lemma}

\begin{proof}
$D_{kl}=0$ if and only if the two non-zero vectors are proportional, which is an equivalence relation; call its classes \emph{lines}. Note that the condition concerns two disjoint non-proportional pairs, and relabelling within the quadruple is free. If one line contains all items but at most one, every pair of distinct items that avoids the exceptional item is proportional. Two disjoint pairs cannot both contain the exceptional item, so no quadruple works. Conversely, suppose no line contains $p-1$ items. If two lines each contain at least two items, say $k_1,k_2$ on one and $k_3,k_4$ on the other, the pairs $(k_1,k_3)$ and $(k_2,k_4)$ are non-proportional. Otherwise every line but at most one is a singleton, and since no line has $p-1\ge3$ items there are at least three lines. Pick $x,y,z$ on three different lines and a fourth item $u$. If $u$ is not on the line of $x$, the pairs $(x,u)$ and $(y,z)$ work; otherwise $(u,y)$ and $(x,z)$ do.
\end{proof}

\subsection{Proof of \texorpdfstring{\cref{lem:id-reduction}}{the reduction lemma}}\label{app:pm-reduction}

Under the normal-ogive link, \eqref{eq:model-lr} with $\varepsilon_j\sim\mathrm N(0,1)$ gives $\bm Y^*\sim\mathrm N_J(\bm A\bm\mu,\bm A\bm\Sigma\bm A^\top+\bm I_J)$, and $\bm Y$ is its discretisation at the thresholds. Let $Z_j=(Y^*_j-\bm a_j^\top\bm\mu)/s_j$. Then $\bm Z\sim\mathrm N_J(\bm0,\bm R)$ and $Y_j=\sum_c\ind{Z_j>\bar b_{jc}}$, so the law of $\bm Y$ depends on the parameters only through $(\bar{\bm b},\bm R)$. Conversely, $\Pr(Y_j\ge c)=\Phi(-\bar b_{jc})$ determines $\bar b_{jc}$. For $j\neq k$, $\Pr(Y_j\ge1,Y_k\ge1)=\Phi_2(-\bar b_{j1},-\bar b_{k1};R_{jk})$, where $\Phi_2(\cdot,\cdot;\rho)$ is the standard bivariate normal distribution function with correlation $\rho$. Since $\partial\Phi_2(x,y;\rho)/\partial\rho=\phi_2(x,y;\rho)>0$ for $|\rho|<1$ \citep{plackett1954}, and $|R_{jk}|<1$ because the error variances are positive, this probability determines $R_{jk}$. In the reference group, $s_j^2=1+\|\bm a_j\|^2$ and $h_j^2=\|\bm a_j\|^2/(1+\|\bm a_j\|^2)$, so $s_j^{-2}=1-h_j^2$, $R_{jk}=\bm a_j^\top\bm a_k/(s_js_k)=\bm\lambda_j^\top\bm\lambda_k$ and $\bar b_{jc}=b_{jc}(1-h_j^2)^{1/2}$. Given $\bm\lambda_j$ with $\|\bm\lambda_j\|^2<1$, the map $\bm a_j=\bm\lambda_j(1-\|\bm\lambda_j\|^2)^{-1/2}$ satisfies $1+\|\bm a_j\|^2=(1-\|\bm\lambda_j\|^2)^{-1}$ and hence returns $\bm\lambda_j$. This establishes the bijection and its inverse. \qed

\subsection{Proof of \texorpdfstring{\cref{thm:ident}}{the identification theorem}}\label{app:pm-ident}

By \cref{lem:id-reduction}, two reference-group parameter values are observationally equivalent if and only if they share $\bar{\bm b}$ and $\operatorname{offdiag}(\bm\Lambda\bm\Lambda^\top)$. Because the map to $(\bm A,\bm b)$ is a bijection, it suffices to determine the set of $\tilde{\bm\Lambda}$ with the pattern \eqref{eq:model-pattern}, communalities below one, and $\tilde{\bm\Lambda}\tilde{\bm\Lambda}^\top$ off-diagonally equal to $\bm\Lambda\bm\Lambda^\top$. Fix such a $\tilde{\bm\Lambda}$.

\paragraph{Step 1: the cross-domain block}
For $k\in\mathcal S$ and $j\in\mathcal I$, $(\bm\Lambda\bm\Lambda^\top)_{kj}=\lambda_{kG}\lambda_{jG}$, because the specific factors do not overlap. Hence $\tilde{\bm w}^{\mathcal S}(\tilde{\bm w}^{\mathcal I})^\top=\bm w^{\mathcal S}(\bm w^{\mathcal I})^\top$. Under \cref{asm:id-gen} this matrix is non-zero, and \cref{lem:app-rank1} gives $t\neq0$ with $\tilde{\bm w}^{\mathcal S}=t\bm w^{\mathcal S}$ and $\tilde{\bm w}^{\mathcal I}=\bm w^{\mathcal I}/t$. Reflecting $G$ if necessary, we may take $t>0$.

\paragraph{Step 2: within a domain}
Within $\mathcal S$, $\tilde{\bm w}\tilde{\bm w}^\top+\tilde{\bm s}\tilde{\bm s}^\top$ is off-diagonally equal to $\bm w\bm w^\top+\bm s\bm s^\top$. By Step 1 this says that $\tilde{\bm s}\tilde{\bm s}^\top$ is off-diagonally equal to $\bm s\bm s^\top+\delta_{\mathcal S}\bm w\bm w^\top$ with $\delta_{\mathcal S}=1-t^2$. Within $\mathcal I$ the same holds with $\delta_{\mathcal I}=1-t^{-2}$.

\paragraph{Part (a)}
Let $\mathcal D$ be the domain in \cref{asm:id-het}. By \cref{lem:app-tetrad}, the tetrad of $\tilde{\bm s}\tilde{\bm s}^\top$ on $(k_1,\dots,k_4)$ is zero and equals $\delta_{\mathcal D}D_{k_1k_4}D_{k_2k_3}$. Hence $\delta_{\mathcal D}=0$, and since $t>0$ both conditions $\delta_{\mathcal S}=0$ and $\delta_{\mathcal I}=0$ are equivalent to $t=1$. So $\tilde{\bm w}=\bm w$ in both domains, and in each domain $\tilde{\bm s}\tilde{\bm s}^\top$ is off-diagonally equal to $\bm s\bm s^\top$. By \cref{asm:id-spec} and \cref{lem:app-onefactor}, $\tilde{\bm s}=\pm\bm s$, which is a reflection of the specific factor. Thus $\tilde{\bm\Lambda}$ equals $\bm\Lambda$ up to column signs, and so do the item parameters, by \cref{lem:id-reduction}. The ECVs \eqref{eq:model-ecv} depend on squared loadings only, so they and the sets $\mathcal B_\tau$ are identified.

\paragraph{Part (b): construction of the curve}
We first show that in each domain $\mathcal D$ satisfying \cref{asm:id-spec} but not \cref{asm:id-het} there is a continuous map $\delta\mapsto\bm u(\delta)$ on a neighbourhood of $0$ with $\bm u(0)=\bm s$, such that $\bm u(\delta)\bm u(\delta)^\top$ is off-diagonally equal to $\bm M(\delta)=\bm s\bm s^\top+\delta\bm w\bm w^\top$.
\begin{itemize}
  \item If $|\mathcal D|=3$, \cref{asm:id-spec} makes all three $s_k$ non-zero. For $\{k,l,m\}=\mathcal D$ put $u_k(\delta)=\operatorname{sgn}(s_k)\{M_{kl}(\delta)M_{km}(\delta)/M_{lm}(\delta)\}^{1/2}$. For small $|\delta|$ each $M_{kl}(\delta)$ has the sign of $s_ks_l\neq0$, so the ratio is positive and continuous, and it equals $s_k^2$ at $\delta=0$. Moreover $u_ku_l=\operatorname{sgn}(s_ks_l)\,|M_{kl}|=M_{kl}$.
  \item If $|\mathcal D|\ge4$, \cref{lem:app-classes} gives a unit vector $(w^*,s^*)$ and scalars $\alpha_k$ with $(w_k,s_k)=\alpha_k(w^*,s^*)$ for all $k\in\mathcal D$ except at most one item $o$. Line items have $s_k=\alpha_ks^*$, so if $s^*=0$ at most one item of $\mathcal D$ would have a non-zero specific loading, contradicting \cref{asm:id-spec}. Orient the line so that $s^*>0$, put $r(\delta)=(s^{*2}+\delta w^{*2})^{1/2}$, which is positive near $0$, and set $u_k=\alpha_kr(\delta)$ on the line and $u_o=(s^*s_o+\delta w^*w_o)/r(\delta)$. Then $u_ku_l=\alpha_k\alpha_l(s^{*2}+\delta w^{*2})=M_{kl}$ and $u_ku_o=\alpha_k(s^*s_o+\delta w^*w_o)=M_{ko}$, and $\bm u(0)=\bm s$.
\end{itemize}
Now let both domains fail \cref{asm:id-het}. For $t$ near $1$ define $\tilde{\bm w}^{\mathcal S}=t\bm w^{\mathcal S}$, $\tilde{\bm w}^{\mathcal I}=\bm w^{\mathcal I}/t$, $\tilde{\bm s}^{\mathcal S}=\bm u^{\mathcal S}(1-t^2)$ and $\tilde{\bm s}^{\mathcal I}=\bm u^{\mathcal I}(1-t^{-2})$. The cross-domain block is $\bm w^{\mathcal S}(\bm w^{\mathcal I})^\top$, as required. Within $\mathcal S$, $\tilde{\bm w}\tilde{\bm w}^\top+\tilde{\bm s}\tilde{\bm s}^\top$ is off-diagonally equal to $t^2\bm w\bm w^\top+\bm s\bm s^\top+(1-t^2)\bm w\bm w^\top=\bm w\bm w^\top+\bm s\bm s^\top$, and similarly within $\mathcal I$. All communalities are continuous in $t$ and below one at $t=1$, so they stay below one on an open interval $T\ni1$. Mapping back with \cref{lem:id-reduction}, keeping $\bar{\bm b}$ fixed, gives a continuous curve of observationally equivalent $(\bm A(t),\bm b(t))$ satisfying \eqref{eq:id-curve}. It is injective because $\bm w^{\mathcal I}\neq\bm0$.

\paragraph{Part (b): the ECVs move}
Let $\epsilon=1-t^{-2}$, which is strictly increasing in $t$, and suppose all loadings in $\mathcal I$ are non-zero. For $j\in\mathcal I$ the odds $\mathrm{ECV}_j/(1-\mathrm{ECV}_j)$ along the curve equal $O_j(\epsilon)=w_j^2(1-\epsilon)/u_j(\epsilon)^2$.
\begin{itemize}
  \item If $|\mathcal I|=3$, put $x_k=w_k/s_k\neq0$. Then $u_j^2=s_j^2(1+\epsilon x_jx_l)(1+\epsilon x_jx_m)/(1+\epsilon x_lx_m)$, and
  \[
    O_j(\epsilon)=x_j^2\,\frac{(1-\epsilon)(1+\epsilon x_lx_m)}{(1+\epsilon x_jx_l)(1+\epsilon x_jx_m)} .
  \]
  If $O_j$ were constant on an interval, the two quadratic polynomials $(1-\epsilon)(1+\epsilon x_lx_m)$ and $(1+\epsilon x_jx_l)(1+\epsilon x_jx_m)$ would coincide. Their $\epsilon^2$ coefficients are $-x_lx_m$ and $x_j^2x_lx_m$, so $x_lx_m(1+x_j^2)=0$, which is impossible. Hence every $\mathrm{ECV}_j$ is non-constant.
  \item If $|\mathcal I|\ge4$, every item on the common line has $O_j(\epsilon)=w^{*2}(1-\epsilon)/(s^{*2}+\epsilon w^{*2})$, where $w^*\neq0$ because $w_j\neq0$. Its derivative is $-(s^{*2}+w^{*2})/(s^{*2}+\epsilon w^{*2})^2<0$, so the odds are strictly decreasing. At least $|\mathcal I|-1$ items lie on the line, and all of them do when the loadings are proportional.
\end{itemize}
A continuous non-constant function of $t$ on $T$ has a range with non-empty interior. For $\tau$ in that interior, some equivalent parameter values place $j$ in $\mathcal B_\tau$ and others do not.

\paragraph{Part (c)}
Steps 1--2 and the argument of part (a), applied to $\mathcal S$, give $t=1$. Hence all general loadings are identified, and $\bm\lambda_{\mathcal SS}$ is identified up to sign by \cref{lem:app-onefactor}. The Israel domain contributes the single off-diagonal entry $R_{j_1j_2}=\lambda_{j_1G}\lambda_{j_2G}+\lambda_{j_1I}\lambda_{j_2I}$, which identifies $\kappa$. Every $(\tilde\lambda_{j_1I},\tilde\lambda_{j_2I})=(q,\kappa/q)$ that keeps both communalities below one is observationally equivalent, and nothing else is. The admissible set is $q^2\in\bigl(\kappa^2/(1-\lambda_{j_2G}^2),\,1-\lambda_{j_1G}^2\bigr)$, which is non-empty because it contains the true value. Since $\mathrm{ECV}_{j_1}=\lambda_{j_1G}^2/(\lambda_{j_1G}^2+q^2)$ is continuous and strictly decreasing in $q^2$, its range is the open interval \eqref{eq:id-two}. If instead $\mathcal S$ fails the condition of \cref{asm:id-het}, the construction of part (b) applies to $\mathcal S$. In $\mathcal I$, the single equation becomes $\tilde\lambda_{j_1I}\tilde\lambda_{j_2I}=R_{j_1j_2}-\lambda_{j_1G}\lambda_{j_2G}/t^2$, which is solvable for $t$ near one. The general loadings of $j_1,j_2$, which equal $\lambda_{jG}/t$, then vary together with their specific loadings.

\paragraph{Part (d)}
The correlated-factors model is the normal-ogive model with $\bm\Sigma=\left(\begin{smallmatrix}1&\rho\\ \rho&1\end{smallmatrix}\right)$, so \cref{lem:id-reduction} applies. Write $\tilde\alpha_j$ for the correlation of $Y^*_j$ with its factor. Then $R_{kl}=\tilde\alpha_k\tilde\alpha_l$ within a domain and $R_{kj}=\rho\tilde\alpha_k\tilde\alpha_j$ across domains. For $k,l\in\mathcal S$ and $j,m\in\mathcal I$ with non-zero loadings, $R_{kj}R_{lm}/(R_{kl}R_{jm})=\rho^2$. If $\rho\neq0$, its sign is that of $R_{k_0j_0}$ for the positively loading marker items $k_0,j_0$. With at least three non-zero loadings per factor, \cref{lem:app-onefactor} gives the loadings up to sign, and the markers fix the signs. With two, say $\mathcal S\supseteq\{k,l\}$, and $\rho\neq0$, we have $\tilde\alpha_k^2=R_{kl}R_{kj}/R_{lj}$. \qed

\subsection{Proof of the claims in \texorpdfstring{\cref{rem:id-null}}{the remark on a null general factor}}\label{app:pm-null}

Let $\bm w^{\mathcal I}=\bm0$. Suppose first that $\mathcal S$ satisfies the condition of \cref{asm:id-het}, and let $\tilde{\bm\Lambda}$ be equivalent. The cross-domain block vanishes, so $\tilde{\bm w}^{\mathcal S}=\bm0$ or $\tilde{\bm w}^{\mathcal I}=\bm0$. If $\tilde{\bm w}^{\mathcal S}=\bm0$, then $\bm w\bm w^\top+\bm s\bm s^\top$ on $\mathcal S$ would be off-diagonally equal to a rank-one matrix. Its tetrad would then vanish, but by \cref{lem:app-tetrad} with $\delta=1$ the tetrad is $D_{k_1k_4}D_{k_2k_3}\neq0$. Hence $\tilde{\bm w}^{\mathcal I}=\bm0$ and $\mathrm{ECV}_j=0$ for all $j\in\mathcal I$. On $\mathcal S$, replacing $[\bm w^{\mathcal S},\bm s^{\mathcal S}]$ by $[\bm w^{\mathcal S},\bm s^{\mathcal S}]\bm Q$ for any $2\times2$ orthogonal $\bm Q$ leaves $\bm w\bm w^\top+\bm s\bm s^\top$ unchanged, so the split between $G$ and $S$ is determined only up to rotation.

Suppose instead that $\bm w\bm w^\top+\bm s\bm s^\top$ on $\mathcal S$ is off-diagonally equal to $\bm u\bm u^\top$ with all $u_k^2<1$. This holds for proportional loadings, with $u_k=\pm h_k$. Then for every $\varphi\in[0,\pi/2)$ the parameter value with $\tilde{\bm w}^{\mathcal S}=\bm0$, $\tilde{\bm s}^{\mathcal S}=\bm u$, $\tilde{\bm w}^{\mathcal I}=\sin\varphi\,\bm s^{\mathcal I}$ and $\tilde{\bm s}^{\mathcal I}=\cos\varphi\,\bm s^{\mathcal I}$ is observationally equivalent, with $\mathrm{ECV}_j=\sin^2\varphi$ for every $j\in\mathcal I$. Finally, along the smooth curve $\bm Q(\varphi)$ of rotations the law of $\bm Y$ is constant, so the score in the tangent direction vanishes identically, and that direction lies in the null space of the Fisher information. \qed

\subsection{Proof of \texorpdfstring{\cref{prop:model-sl}}{the Schmid--Leiman proposition}}\label{app:pm-sl}

\paragraph{Part (a)}
Put $\tilde\zeta_D=\gamma_D\theta_G+(1-\gamma_D^2)^{1/2}\theta_D$ for $D\in\{S,I\}$. With $\theta_G,\theta_S,\theta_I$ independent standard normal, $(\tilde\zeta_S,\tilde\zeta_I)$ is bivariate normal with unit variances and covariance $\gamma_S\gamma_I=\rho_{SI}$, so it has the law of $(\zeta_S,\zeta_I)$. Under \eqref{eq:model-sl} the linear predictor of item $j$ is $a_{jG}\theta_G+a_{js}\theta_{s(j)}=\alpha_j\tilde\zeta_{s(j)}$. Hence the vector of linear predictors has the same law as under the correlated-factors model. Since the errors are independent of it with the same distribution $F$, the latent responses, and therefore the responses, have the same law for every link. Finally, $\mathrm{ECV}_j=\alpha_j^2\gamma_{s(j)}^2/\{\alpha_j^2\gamma_{s(j)}^2+\alpha_j^2(1-\gamma_{s(j)}^2)\}=\gamma_{s(j)}^2$ for $\alpha_j\neq0$.

\paragraph{Part (b)}
If the loading vectors of domain $D$ lie on the line spanned by $(\cos\phi_D,\sin\phi_D)$, write $(a_{jG},a_{js})=\alpha_j(\cos\phi_D,\sin\phi_D)$. After reflecting the specific factor if necessary, $\sin\phi_D\ge0$. Then $\gamma_D=\cos\phi_D$ gives \eqref{eq:model-sl} with $\rho_{SI}=\gamma_S\gamma_I$.

For the converse under the normal-ogive link, let the bifactor model with standardised loadings $(\bm w,\bm s)$ be observationally equivalent to a correlated-factors model with standardised loadings $\tilde{\bm\alpha}$ and correlation $\rho$. By \cref{lem:id-reduction} the two polychoric matrices are off-diagonally equal. Across domains, $\bm w^{\mathcal S}(\bm w^{\mathcal I})^\top=\rho\,\tilde{\bm\alpha}^{\mathcal S}(\tilde{\bm\alpha}^{\mathcal I})^\top$. The left side is non-zero, so $\rho\neq0$, and \cref{lem:app-rank1} gives $\bm w^{\mathcal S}=c\,\tilde{\bm\alpha}^{\mathcal S}$ for some $c\neq0$. Within $\mathcal S$, $\bm s\bm s^\top$ is then off-diagonally equal to $(1-c^2)\tilde{\bm\alpha}\tilde{\bm\alpha}^\top$. Take three items $k,l,m$ with non-zero $s$. Then $(1-c^2)\tilde\alpha_k\tilde\alpha_l=s_ks_l\neq0$, and so on for the other pairs, which gives $s_k^2=(s_ks_l)(s_ks_m)/(s_ls_m)=(1-c^2)\tilde\alpha_k^2$. Hence $1-c^2>0$. Applying \cref{lem:app-onefactor} to $\bm u=(1-c^2)^{1/2}\tilde{\bm\alpha}$, which has at least three non-zero entries, yields $\bm s=\pm(1-c^2)^{1/2}\tilde{\bm\alpha}$. Therefore $(w_k,s_k)=\tilde\alpha_k\bigl(c,\pm(1-c^2)^{1/2}\bigr)$ for all $k\in\mathcal S$, which are proportional vectors. The same argument applies to $\mathcal I$.

\paragraph{Part (c)}
By (a), every point of $\{\gamma_S\gamma_I=\rho_{SI},\ |\gamma_S|,|\gamma_I|\le1\}$ is a bifactor representation. Under the conditions of (b), any representation has proportional loadings and is therefore of this form, with $\rho$ equal to $\rho_{SI}$ by \cref{thm:ident}(d). On the curve $\mathrm{ECV}_j=\gamma_I^2$ for $j\in\mathcal I$, and $|\gamma_S|=|\rho_{SI}|/|\gamma_I|\le1$ forces $\gamma_I^2\ge\rho_{SI}^2$. Every $v\in[\rho_{SI}^2,1]$ is attained, with $\gamma_I=v^{1/2}$ and $\gamma_S=\rho_{SI}v^{-1/2}$; for $v=0$, which requires $\rho_{SI}=0$, any $\gamma_S$ works. At $v=1$ all $I$-loadings vanish, and at $v=\rho_{SI}^2$ all $S$-loadings do. For $k\in\mathcal S$, $\mathrm{ECV}_k=\gamma_S^2=\rho_{SI}^2/\gamma_I^2$. \qed

\subsection{Proof of \texorpdfstring{\cref{thm:dif-bias}}{the DIF-bias theorem}}\label{app:pm-dif}

\begin{lemma}[First-order pseudo-true values]\label{lem:app-pseudo}
Let $\mathcal Y$ be finite, and let $\{p(\cdot\,;\bm x):\bm x\in U\}$, with $U\subseteq\mathbb R^k$ open and $\bm0\in U$, be a family of positive probability mass functions on $\mathcal Y$, three times continuously differentiable in $\bm x$. Let $\bm H=\sum_y p(y;\bm0)\bm\sigma(y)\bm\sigma(y)^\top$, where $\bm\sigma(y)=\nabla_{\bm x}\log p(y;\bm0)$. Let $\bm x(\cdot)$ be a $C^3$ map from an open set containing a compact $\Xi_c\subset\mathbb R^q$ into $U$, with $\bm x(\bm\xi_0)=\bm0$ for some $\bm\xi_0$ in the interior of $\Xi_c$ and Jacobian $\bm L$ at $\bm\xi_0$. Let $K(\bm\xi;\bm x)=\sum_yp(y;\bm x)\log\{p(y;\bm x)/p(y;\bm x(\bm\xi))\}$, and suppose that (i) $\bm\xi_0$ is the unique minimiser of $K(\cdot\,;\bm0)$ on $\Xi_c$, and (ii) $\bm L^\top\bm H\bm L$ is non-singular. Then there is a neighbourhood $N$ of $\bm0$ such that, for every $\bm x_\Delta\in N$, $K(\cdot\,;\bm x_\Delta)$ has a unique minimiser $\bm\xi^*(\bm x_\Delta)$ on $\Xi_c$. It lies in the interior of $\Xi_c$, and
\[
  \bm\xi^*(\bm x_\Delta)=\bm\xi_0+(\bm L^\top\bm H\bm L)^{-1}\bm L^\top\bm H\,\bm x_\Delta+O(\|\bm x_\Delta\|^2).
\]
\end{lemma}

\begin{proof}
$K$ is jointly continuous on $\Xi_c\times\bar N_0$ for a compact neighbourhood $\bar N_0$ of $\bm0$, because it is a finite sum of continuous terms with positive arguments. Minimisers exist by compactness. Let $\bm x_n\to\bm0$ and let $\bm\xi_n$ minimise $K(\cdot\,;\bm x_n)$. Since $K(\bm\xi_n;\bm x_n)\le K(\bm\xi_0;\bm x_n)\to K(\bm\xi_0;\bm0)=0$, every limit point $\bar{\bm\xi}$ satisfies $K(\bar{\bm\xi};\bm0)\le0$, so $\bar{\bm\xi}=\bm\xi_0$ by (i). Hence minimisers converge to $\bm\xi_0$; for small $\bm x$ they are interior and solve $\bm\Gamma(\bm\xi,\bm x)=\bm0$, where $\bm\Gamma(\bm\xi,\bm x)=\nabla_{\bm\xi}K=-\bm L(\bm\xi)^\top\sum_yp(y;\bm x)\nabla_{\bm x}\log p(y;\bm x(\bm\xi))$ and $\bm L(\bm\xi)=\partial\bm x/\partial\bm\xi^\top$. At $(\bm\xi_0,\bm0)$ we have $\sum_yp\,\bm\sigma=\bm0$ and $-\sum_yp\,\nabla^2_{\bm x}\log p=\bm H$, so $\partial\bm\Gamma/\partial\bm\xi^\top=\bm L^\top\bm H\bm L$. The weights depend on $\bm x$ through $\nabla_{\bm x}p=p\,\bm\sigma$, so $\partial\bm\Gamma/\partial\bm x^\top=-\bm L^\top\sum_yp\,\bm\sigma\bm\sigma^\top=-\bm L^\top\bm H$. $\bm\Gamma$ is $C^2$ and its $\bm\xi$-derivative is invertible by (ii). The implicit function theorem gives a unique $C^2$ solution $\bm\xi(\bm x)$ near $(\bm\xi_0,\bm0)$ with derivative $(\bm L^\top\bm H\bm L)^{-1}\bm L^\top\bm H$ at $\bm0$. Every minimiser for small $\bm x$ lies in this neighbourhood and solves the equation, so it equals $\bm\xi(\bm x)$. Taylor's theorem gives the expansion.
\end{proof}

\paragraph{The two exact facts}
Let $\bm D=\operatorname{diag}(\sigma_G,\sigma_S,\sigma_I)$ and $\bm z=\bm D^{-1}(\bm\theta-\bm\mu)\sim\mathrm N_3(\bm0,\bm I)$. Under the invariant model with focal parameters $(\bm\mu',\bm\ell')$, write $\bm\theta'=\bm\mu'+\bm D'\bm z$. The predictor of item $j$ is then
\[
  \bm a_j^\top\bm\theta'-b_{jc}=\bm a_j^\top\bm\theta-b_{jc}+\bm a_j^\top(\bm\mu'-\bm\mu)+\sum_f a_{jf}\bigl(\sigma'_f/\sigma_f-1\bigr)(\theta_f-\mu_f).
\]
This is \eqref{eq:id-coords} with $\bm\iota=\bm A(\bm\mu'-\bm\mu)$ and $\beta_{jf}=a_{jf}(e^{(\ell'_f-\ell_f)/2}-1)$. Only $f\in\{G,s(j)\}$ contribute, by the pattern \eqref{eq:model-pattern}, and the Jacobian at $\bm\xi'=\bm\xi$ is $\operatorname{diag}(\bm A,\bm B)$. The identity for $\bm x_\Delta$ was shown in the text; it is exact and linear in $\Delta$.

\paragraph{Part (a)}
Identification of $\bm\xi$ in the invariant model is hypothesis (i) of \cref{lem:app-pseudo}, and non-singularity of $\bm L^\top\bm H\bm L$ is (ii). The lemma, with $\bm x_\Delta$ linear in $\Delta$, gives $\bm\xi^*-\bm\xi=(\bm L^\top\bm H\bm L)^{-1}\bm L^\top\bm H\bm x_\Delta+O(\|\Delta\|^2)$. With $\bm L=\operatorname{diag}(\bm A,\bm B)$, the first-order term $(\Delta\bm\mu,\Delta\bm\ell)$ solves
\[
  \begin{pmatrix}\bm P&\bm C\\ \bm C^\top&\bm Q\end{pmatrix}\begin{pmatrix}\Delta\bm\mu\\ \Delta\bm\ell\end{pmatrix}
  =\begin{pmatrix}\bm A^\top(\bm H_{\iota\iota}\bm d+\bm H_{\iota\beta}\bm\delta)\\ \bm B^\top(\bm H_{\beta\iota}\bm d+\bm H_{\beta\beta}\bm\delta)\end{pmatrix},
\]
with $\bm P=\bm A^\top\bm H_{\iota\iota}\bm A$, $\bm C=\bm A^\top\bm H_{\iota\beta}\bm B$ and $\bm Q=\bm B^\top\bm H_{\beta\beta}\bm B$. Eliminating $\Delta\bm\ell=\bm Q^{-1}\{\bm B^\top(\bm H_{\beta\iota}\bm d+\bm H_{\beta\beta}\bm\delta)-\bm C^\top\Delta\bm\mu\}$ gives $(\bm P-\bm C\bm Q^{-1}\bm C^\top)\Delta\bm\mu=\bm A^\top\bm W_\iota\bm d+\bm A^\top\bm H_{\iota\beta}(\bm I-\Pi_{\bm B})\bm\delta$, and $\bm P-\bm C\bm Q^{-1}\bm C^\top=\bm A^\top\bm W_\iota\bm A$. This is \eqref{eq:id-dif-mu}. The symmetric elimination gives \eqref{eq:id-dif-ell}. Both Schur complements are non-singular because $\bm L^\top\bm H\bm L$ is.

\paragraph{Part (b): symmetry}
Write $\bar\tau_{jc}=b_{jc}-\bm a_j^\top\bm\mu$, so the hypothesis is $\bar\tau_{j,C_j+1-c}=-\bar\tau_{jc}$; with $\bar\tau_{j0}=-\infty$ and $\bar\tau_{j,C_j+1}=+\infty$ it holds for $c=0,\dots,C_j+1$. Let $\bm z=\bm\theta-\bm\mu$, whose law is symmetric, and let $\omega_j=\bm a_j^\top\bm z+\iota_j+\bm\beta_j^\top\bm z$ be the perturbed predictor, so that $\Pr(Y_j=k\mid\bm z)=F(\omega_j-\bar\tau_{jk})-F(\omega_j-\bar\tau_{j,k+1})=:\pi_k(\omega_j)$. By symmetry of $F$,
\[
  \pi_{C_j-k}(\omega)=F(\omega+\bar\tau_{j,k+1})-F(\omega+\bar\tau_{jk})=F(-\omega-\bar\tau_{jk})-F(-\omega-\bar\tau_{j,k+1})=\pi_k(-\omega).
\]
Let $\mathsf R\bm y=(C_j-y_j)_j$. Then $p_1(\mathsf R\bm y;\bm\iota,\bm\beta)=\E\prod_j\pi_{y_j}(-\omega_j)$. Since $-\omega_j=\bm a_j^\top(-\bm z)-\iota_j+\bm\beta_j^\top(-\bm z)$ and $-\bm z$ has the law of $\bm z$, this equals $p_1(\bm y;-\bm\iota,\bm\beta)$. At $\bm x=\bm0$, therefore, $p_1(\mathsf R\bm y)=p_1(\bm y)$, $\nabla_\iota\log p_1(\mathsf R\bm y)=-\nabla_\iota\log p_1(\bm y)$ and $\nabla_\beta\log p_1(\mathsf R\bm y)=\nabla_\beta\log p_1(\bm y)$. Because $\mathsf R$ is a bijection of the pattern space, $\bm H_{\iota\beta}=-\bm H_{\iota\beta}=\bm0$. With $\bm H_{\iota\beta}=\bm0$, $\bm W_\iota=\bm H_{\iota\iota}$ and $\bm W_\beta=\bm H_{\beta\beta}$, and the formulas reduce as stated. The remaining claims are read off: $\bm d$ enters only the $\bm\mu$-equation, and $\bm d=\Delta\bm A\bm\mu$ when $\Delta\bm b=\bm0$.

\paragraph{Part (b): linear-Gaussian case}
Suppose $\bm Y^*$ is observed. The truth is $\mathrm N_J(\bm A\bm\mu+\bm d,\tilde{\bm V})$ and the model is $\mathrm N_J(\bm A\bm\mu',\bm V')$ with $\bm V'=\bm A\bm\Sigma'\bm A^\top+v_F\bm I$. The Kullback--Leibler divergence equals $\tfrac12\{\operatorname{tr}(\bm V'^{-1}\tilde{\bm V})+(\bm A\bm\mu+\bm d-\bm A\bm\mu')^\top\bm V'^{-1}(\cdot)+\log\det\bm V'-\log\det\tilde{\bm V}-J\}$. For fixed $\bm\Sigma'\succ0$ it is minimised over $\bm\mu'$ by generalised least squares, $\bm\mu'=\bm\mu+(\bm A^\top\bm V'^{-1}\bm A)^{-1}\bm A^\top\bm V'^{-1}\bm d$. Let $\bm N=\bm A^\top\bm A/v_F$ and $\bm K=(\bm\Sigma'^{-1}+\bm N)^{-1}$. By the Woodbury identity $\bm V'^{-1}=v_F^{-1}\bm I-v_F^{-2}\bm A\bm K\bm A^\top$. Hence $\bm A^\top\bm V'^{-1}=(\bm I-\bm N\bm K)\bm A^\top/v_F=\bm\Sigma'^{-1}\bm K\bm A^\top/v_F$ and $\bm A^\top\bm V'^{-1}\bm A=\bm\Sigma'^{-1}\bm K\bm N$. Therefore $(\bm A^\top\bm V'^{-1}\bm A)^{-1}\bm A^\top\bm V'^{-1}=\bm N^{-1}\bm A^\top/v_F=(\bm A^\top\bm A)^{-1}\bm A^\top$. This does not depend on $\bm\Sigma'$, so the joint minimiser has $\bm\mu^*=\bm\mu+(\bm A^\top\bm A)^{-1}\bm A^\top\bm d$ for every $\Delta$.

\paragraph{Part (c)}
Apply \eqref{eq:model-logit-ecv} to $(\bm a_j,\bm\ell^*)$ and to $(\bm a_j+\Delta\bm a_j,\bm\ell)$ and subtract. The identity is exact whenever $a_{jG}+\Delta a_{jG}$ and $a_{js}+\Delta a_{js}$ are non-zero.

\paragraph{Part (d)}
(i) $\bm a_j^\top\bm\theta-b_{jc}-\Delta b_j=\bm a_j^\top(\bm\theta+\bm v)-b_{jc}$ when $\Delta b_j=-\bm a_j^\top\bm v$, and $\bm\theta+\bm v\sim\mathrm N(\bm\mu+\bm v,\bm\Sigma)$. The invariant model at $(\bm\mu+\bm v,\bm\ell)$ thus reproduces the law, the divergence is zero, and by identification $\bm\xi^*=(\bm\mu+\bm v,\bm\ell)$.
(ii) If $-\bm A\bm v$ vanishes on $\mathcal S$, then $v_Ga_{kG}+v_Sa_{kS}=0$ for all $k\in\mathcal S$. Since the stereotype loading vectors span $\mathbb R^2$, $v_G=v_S=0$ and $\Delta\bm b_{\mathcal I}=-v_I\bm a_{\mathcal II}$. By (i) such DIF changes only $\mu_I$, so $\bm\ell$ and all ECVs are unaffected. In the linear-Gaussian case the bias of $\mu_G$ for DIF supported on $\mathcal I$ is $\bm e_G^\top(\bm A^\top\bm A)^{-1}\sum_{j\in\mathcal I}\bm a_jd_j$. It vanishes for all such $\bm d$ only if $\bm e_G^\top(\bm A^\top\bm A)^{-1}$ is orthogonal to $\{(a_{jG},0,a_{jI}):j\in\mathcal I\}$. If these vectors are not all proportional they span $\{\bm e_G,\bm e_I\}$, and orthogonality would give $\bm e_G^\top(\bm A^\top\bm A)^{-1}\bm e_G=0$, which contradicts positive definiteness.
(iii) This is \eqref{eq:id-dif-ecv}. \qed

\subsection{Proof of \texorpdfstring{\cref{prop:anchor}}{the anchor proposition}}\label{app:pm-anchor}

\paragraph{Part (a)}
Let $m_j=\bm a_j^\top\bm\mu_1$ and $\tilde s_j=(1+\bm a_j^\top\bm\Sigma_1\bm a_j)^{1/2}$ for an anchor $j$. By \cref{lem:id-reduction}, $\Pr(Y_j\ge c\mid g=1)=\Phi\{(m_j-b_{jc})/\tilde s_j\}$ for $c=1,2$, with $b_{j1}<b_{j2}$ known. Writing $z_c$ for the two normal quantiles, $\tilde s_j=(b_{j2}-b_{j1})/(z_1-z_2)$ and $m_j=b_{j1}+\tilde s_jz_1$. The bivariate margins identify the group-1 polychoric correlations, hence $\Cov(Y^*_j,Y^*_k)=R^{(1)}_{jk}\tilde s_j\tilde s_k$ for anchors $j\neq k$. Thus $\bm m_{\mathcal A}=\bm A_{\mathcal A}\bm\mu_1$ and $\bm C_{\mathcal A}=\bm A_{\mathcal A}\bm\Sigma_1\bm A_{\mathcal A}^\top+\bm I$ are identified. Full column rank gives $\bm\mu_1=(\bm A_{\mathcal A}^\top\bm A_{\mathcal A})^{-1}\bm A_{\mathcal A}^\top\bm m_{\mathcal A}$ and $\bm\Sigma_1=\bm A_{\mathcal A}^+(\bm C_{\mathcal A}-\bm I)\bm A_{\mathcal A}^{+\top}$.

For a non-anchor $k$ with group-1 parameters $(\bm a_k,b_{kc})$, the correlations with the anchors identify the vector $\bm A_{\mathcal A}\bm\Sigma_1\bm a_k/\tilde s_k$, with entries $R^{(1)}_{jk}\tilde s_j$. Full column rank then identifies $\bm q=\bm a_k/\tilde s_k$. From $\tilde s_k^2=1+\tilde s_k^2\,\bm q^\top\bm\Sigma_1\bm q$ we get $\bm q^\top\bm\Sigma_1\bm q<1$, $\tilde s_k^2=(1-\bm q^\top\bm\Sigma_1\bm q)^{-1}$ and $\bm a_k=\tilde s_k\bm q$. Finally $\Pr(Y_k\ge c\mid g=1)=\Phi\{(\bm a_k^\top\bm\mu_1-b_{kc})/\tilde s_k\}$ identifies $b_{kc}$. Without anchors, $\bm a_j^{(1)\top}\bm\theta-b^{(1)}_{jc}=\bm a_j^{(1)\top}(\bm\theta+\bm v)-(b^{(1)}_{jc}+\bm a_j^{(1)\top}\bm v)$ for every item, which is the stated invariance.

\paragraph{Part (b)}
For anchors, $\bm a_j^\top\bm\theta-b_{jc}-\Delta b_j=\bm a_j^\top(\bm\theta+\bm v)-b_{jc}$. For non-anchors, $\bm a_k^{(1)\top}\bm\theta-b^{(1)}_{kc}=\bm a_k^{(1)\top}(\bm\theta+\bm v)-(b^{(1)}_{kc}+\bm a_k^{(1)\top}\bm v)$. The two descriptions therefore induce the same law, and under (a) the anchored analysis converges to the second one.

\paragraph{Part (c)}
Let $\bm z=\bm\Sigma_1^{-1/2}(\bm\theta-\bm\mu_1)$. Reparametrise each non-anchor item by its predictor as a function of $\bm z$, $\bm c_k^\top\bm z-e_{kc}$. For fixed $(\bm\mu_1,\bm\Sigma_1)$ this is a bijection of its parameters, and it makes the non-anchor parameters $\bm\nu=(\bm c_k,e_{k\cdot})_{k\notin\mathcal A}$ variation-independent of $(\bm\mu,\bm\ell)$. In the coordinates $(\bm\iota_{\mathcal A},\bm\beta_{\mathcal A},\bm\nu)$, the anchors' perturbations \eqref{eq:id-coords} together with $\bm\nu$, the anchored model is $\bm\iota_{\mathcal A}=\bm A_{\mathcal A}(\bm\mu'-\bm\mu)$, $\bm\beta_{\mathcal A}$ a function of $\bm\ell'$ with Jacobian $\bm B_{\mathcal A}$, and $\bm\nu$ free. Hence $\bm L=\operatorname{diag}(\bm A_{\mathcal A},\bm B_{\mathcal A},\bm I)$. The truth is at $(\bm d_{\mathcal A},\bm0,\bm0)$ for uniform anchor DIF, or $(\bm d_{\mathcal A},\bm\delta_{\mathcal A},\bm0)$ in general. By \cref{lem:app-pseudo}, the first-order term solves the normal equations of the $\bm H$-weighted projection. Eliminating the free block $\bm\nu$ (and, for uniform DIF, the $\bm\ell$ block) replaces $\bm H$ by its Schur complement, as in the proof of \cref{thm:dif-bias}(a). This gives \eqref{eq:id-anchor-bias} with $\bm W$ the efficient information of $\bm\iota_{\mathcal A}$, and the stated analogue of \eqref{eq:id-dif-mu}--\eqref{eq:id-dif-ell} for non-uniform DIF. For the bound, let $\bm e=(\bm M^\top\bm M)^{-1}\bm M^\top\bm W^{1/2}\bm d_{\mathcal A}$ be the first-order term, where $\bm W^{1/2}$ is the symmetric square root. Then $\bm M\bm e=\Pi\bm W^{1/2}\bm d_{\mathcal A}$ and $\|\bm M\bm e\|_2\ge\sigma_{\min}(\bm M)\|\bm e\|_2$, which gives the first inequality; the second holds because $\Pi$ is an orthogonal projection. For one contaminated anchor, $\bm d_{\mathcal A}=-\Delta b_j\bm e_j$ and $\|\Pi\bm W^{1/2}\bm e_j\|_2^2=\bm e_j^\top\bm W^{1/2}\Pi\bm W^{1/2}\bm e_j=h_j$. Moreover $\bm W^{1/2}\Pi\bm W^{1/2}\preceq\bm W$, so $h_j\le W_{jj}$. \qed

\subsection{Proof of \texorpdfstring{\cref{prop:id-boundary-test}}{the boundary-test proposition}}\label{app:pm-test}

The response pattern takes finitely many values, and its distribution is a smooth function of $\bm\vartheta$. Under the conditions of \cref{thm:ident}(a), with $\bm\vartheta$ interior and Fisher information $\bm{\mathcal F}$ non-singular, the maximum likelihood estimator satisfies $n^{1/2}(\widehat{\bm\vartheta}-\bm\vartheta)\Rightarrow\mathrm N(\bm0,\bm{\mathcal F}^{-1})$, and the inverse observed information is consistent. The function $g_j$ is smooth near $\bm\vartheta$ when $a_{jG}a_{jI}\neq0$, with gradient $(2/a_{jG},-2/a_{jI})$ in $(a_{jG},a_{jI})$. Hence $n^{1/2}(\widehat g_j-g_j)\Rightarrow\mathrm N(0,v_j)$ with $v_j>0$ and $n\,\widehat{\mathrm{se}}_j^2\to v_j$ in probability.

\paragraph{Part (a)}
Write the statistic as $(\widehat g_j-g_j)/\widehat{\mathrm{se}}_j+(g_j-\logit\tau)/\widehat{\mathrm{se}}_j$. The first term converges to $\mathrm N(0,1)$, and the second tends to $0$, $+\infty$ or $-\infty$ according as $g_j=$, $>$ or $<\logit\tau$, because $\logit$ is increasing. The stated limits follow, and the supremum over the null of the limiting rejection probability is $\alpha$. The signed root of the likelihood-ratio statistic for the smooth scalar restriction $g_j=\logit\tau$ differs from the Wald statistic by $o_p(1)$ under the null and local alternatives. For the one-sided null, the likelihood-ratio statistic is zero when $\widehat g_j\ge\logit\tau$; otherwise, with probability tending to one, the restricted maximum lies on $\{g_j=\logit\tau\}$. Its limit is therefore $Z^2\ind{Z<0}$ with $Z\sim\mathrm N(0,1)$, that is $\tfrac12\chi^2_0+\tfrac12\chi^2_1$ \citep{chernoff1954,self1987}.

\paragraph{Part (b)}
On the event $E$ that $L_j\le\mathrm{ECV}_j\le U_j$ for all $j\in\mathcal I$, every $j\in\widehat{\mathcal B}{}^{\mathrm{in}}_\tau$ has $\mathrm{ECV}_j\ge L_j\ge\tau$, and every $j\in\mathcal B_\tau$ has $U_j\ge\mathrm{ECV}_j\ge\tau$. Hence the inclusions hold on $E$ and $\liminf\Pr(E)\ge1-\alpha$. For Bonferroni intervals, each two-sided interval for $g_j$ at level $1-\alpha/|\mathcal I|$ has asymptotic coverage $1-\alpha/|\mathcal I|$, monotone transformation by $\operatorname{expit}$ preserves coverage, and the union bound gives $\liminf\Pr(E)\ge1-\alpha$.

\paragraph{Part (c)}
At $\mathrm{ECV}_j=\tau$, $\Pr(\widehat{\mathrm{ECV}}_j\ge\tau)=\Pr\{n^{1/2}(\widehat g_j-g_j)\ge0\}\to\tfrac12$. At $a_{jI}=0$, $n(1-\widehat{\mathrm{ECV}}_j)=n\widehat a_{jI}^2/(\widehat a_{jG}^2+\widehat a_{jI}^2)$, with $n\widehat a_{jI}^2\Rightarrow V_{jI}\chi^2_1$ and a denominator converging to $a_{jG}^2>0$. Slutsky's theorem gives the limit. \qed
